\documentclass[11pt,a4paper]{article}

\usepackage[margin=2.3cm]{geometry}
\usepackage[T1]{fontenc}
\usepackage{lmodern}
\usepackage{amsmath,amssymb}
\usepackage{booktabs}
\usepackage{tabularx}
\usepackage{algorithm}
\usepackage{algpseudocode}
\usepackage{listings}
\usepackage{xcolor}
\usepackage{enumitem}
\usepackage[colorlinks=true,linkcolor=blue,citecolor=blue,urlcolor=blue]{hyperref}

\definecolor{codegray}{gray}{0.95}
\definecolor{commentgreen}{rgb}{0.25,0.45,0.25}
\lstset{
  basicstyle=\ttfamily\footnotesize,
  backgroundcolor=\color{codegray},
  commentstyle=\color{commentgreen}\itshape,
  keywordstyle=\bfseries,
  breaklines=true,
  columns=fullflexible,
  keepspaces=true,
  frame=none,
  xleftmargin=0.5em,
  showstringspaces=false,
  language=Python,
}
\setlist{leftmargin=*,itemsep=2pt,topsep=3pt}
\newcolumntype{Y}{>{\raggedright\arraybackslash}X}

\newcommand{\zs}{z^{s}}
\newcommand{\ze}{z^{e}}
\newcommand{\zv}{z^{v}}
\newcommand{\xh}{\hat{x}_0}
\newcommand{\xhat}[1]{\hat{x}_{0,#1}}
\newcommand{\sg}{\mathrm{sg}}
\newcommand{\code}[1]{\texttt{#1}}

\title{Latent Recursion in VEDA-S\\\large Design document v4: results, design space, and a proposed training method}
\author{Bjarke L.\,V.\,B.\ Erichsen}
\date{\today}

\begin{document}
\maketitle
\tableofcontents

% ======================================================================
\section{Status}

Working configuration: \code{node\_edge\_e1\_lin\_res}, $K{=}6$, $n{=}6$, with
self-conditioning (SC). It beats SC alone at equal layer count on both QM9 and
GEOM-Drugs. Open: (i) a better training method, (ii) a mechanism for the coordinate
channel. Sections~\ref{sec:mech}--\ref{sec:results} record what exists;
Section~\ref{sec:main} is the main proposal; Section~\ref{sec:coords} the coordinate
channel; Section~\ref{sec:secondary} everything else.

% ======================================================================
\section{Current method}\label{sec:mech}

\paragraph{Mechanism.} Carry node and edge state, not coordinates.
\begin{align*}
  \zs &= \code{pred\_feats}\odot\code{atom\_mask} && \text{after \code{feat\_norm}, } [B,N,384]\\
  \ze &= (\code{pred\_edges}+\code{pred\_edges}^{\top})\odot\code{pair\_mask} && \text{after \code{bond\_norm}, } [B,N,N,128]
\end{align*}
Both are tapped before their readout heads. Injected after \code{feat\_proj} and after
\code{edge\_in\_proj} respectively, both by
\[
  h \leftarrow h + W\,[\,h \,\|\, h - h_{\text{prev}}\,], \qquad W \text{ zero-initialised (with bias).}
\]
Added parameters: $768{\times}384 + 256{\times}128 \approx 0.33$M (a core layer is
roughly 1.8M).

\paragraph{Training.} Algorithm~\ref{alg:current}, as run for the experiments below
(the copy of \code{fm.py} on the remote still has the older version, where $p$ gated
the whole chain). $p$ switches only the untracked warm-up; the $n$ tracked passes
always run. All passes see the same $(x_t,\sigma)$; one loss, on the last pass; the SC
batch $c$ is computed once by a latent-free pass and held fixed for the whole chain.

\begin{algorithm}[h]
\caption{Current training step}\label{alg:current}
\begin{algorithmic}[1]
  \State $c \gets \mathbf{0}$;\quad with prob.\ $\tfrac12$: $c \gets \sg\,f_\theta(x_t,\sigma,\mathbf{0},\varnothing)$ \Comment{self-conditioning pass}
  \State $z \gets \varnothing$
  \State \textbf{if} $u<p$: \textbf{for} $k=1..K$: $z \gets \sg\,f_\theta(x_t,\sigma,c,z)$ \Comment{untracked; the only part $p$ gates}
  \State \textbf{for} $i=1..n$: $z \gets f_\theta(x_t,\sigma,c,z)$ \Comment{tracked, always}
  \State $\hat y \gets f_\theta(x_t,\sigma,c,z)$;\quad \textbf{return} $\mathcal L(\hat y, x_0,\sigma)$
\end{algorithmic}
\end{algorithm}

At $K{=}n{=}6$, $p{=}0.5$, per optimizer step: $0.5 + pK + n + 1 = 10.5$ forward passes
expected, $n+1 = 7$ of them tracked every time, for \emph{one} supervised prediction.
That prediction always has 6 or 12 passes of context behind it. On the cold branch this
is plain end-to-end training of a 7-pass weight-tied network.

% ======================================================================
\section{Results to date}\label{sec:results}

Values read off the W\&B plots; treat as approximate. ``$a \to b$'' is the value at
0.4M and at 1.3M steps.

\begin{center}
\begin{tabular}{lcccc}
\toprule
GEOM-Drugs & validity & fc-validity & mol.\ stability & opt-rmsd \\
\midrule
2 layers, no SC, no LR        & 0.15 & 0.05 & 0.25--0.30 & --- \\
2 layers, SC                  & 0.55--0.60 & $0.30 \to 0.45$ & $0.80 \to 0.87$ & $\sim$1.3 \\
\textbf{2 layers, SC + LR} (0.4M) & \textbf{0.70--0.80} & \textbf{0.50--0.60} & 0.80--0.85 & $\sim$1.2 \\
3 layers, SC (0.35M)          & $\sim$0.70 & & & $\sim$1.25 \\
6 layers, SC (0.2M)           & $\sim$0.85 & & & $\sim$1.12 \\
12 layers, SC (0.13M)         & $\sim$0.85 & $\sim$0.75 & $\sim$0.88 & $\sim$1.03 \\
\bottomrule
\end{tabular}

\smallskip
\begin{tabular}{lcc}
\toprule
QM9 & validity & opt-rmsd \\
\midrule
2 layers, no SC, no LR & 0.83 & 0.68 \\
2 layers, SC           & 0.97--0.99 & 0.39 \\
2 layers, LR only      & 0.97--0.99 & 0.36 \\
2 layers, SC + LR      & 0.97--0.99 & $\sim$0.39 (short run) \\
12 layers, SC          & 0.97--0.99 & 0.27 \\
\bottomrule
\end{tabular}
\end{center}

\begin{itemize}
  \item \textbf{GEOM is where the effect is.} LR at 0.4M steps beats SC at 1.3M steps
        on validity ($\sim$0.75 vs $\sim$0.57), fc-validity (0.6 vs 0.45) and
        energy-validity (0.8 vs 0.55). The LR run is roughly 2.5--5$\times$ slower per
        step, so that is about matched wall-clock: the gain is not only speed. With
        0.33M extra parameters the 2-layer model lands between the 3- and 6-layer SC
        models on validity.
  \item \textbf{Molecule stability is the exception.} SC gets to $\sim$0.87 given
        enough steps; LR sits at 0.80--0.85.
  \item \textbf{QM9 is saturated} on validity (the discrete heads are solved at 2
        layers); only opt-rmsd separates, slightly.
  \item \textbf{Cost.} Per parameter and per \emph{inference} pass LR wins (inference is
        still one pass per step). Per \emph{training} hour the 3- and 6-layer models are
        ahead.
  \item \textbf{Train loss is flattering.} \code{train-type-loss} drops from 0.105 (SC)
        to 0.08 (LR), but it is measured after 6 or 12 same-$\sigma$ refinement passes,
        i.e.\ on a 14- or 26-layer weight-tied network. Inference gets one pass per step.
  \item \textbf{Not usable as they stand.} \code{val-energy} and
        \code{val-energy-per-atom} average over valid molecules only, and LR validates
        $\sim$40\% more of them (use median on the intersection,
        Appendix~\ref{app:metrics}). \code{val-fcd-score} reads a constant $10^6$ for
        most runs, which looks like a failure value. \code{val-novelty} is $\approx 1$
        everywhere.
\end{itemize}

% ======================================================================
\section{Design space}\label{sec:space}

Your table, compressed, with the closest prior work and where each idea stands here.

\begin{center}\small
\begin{tabularx}{\textwidth}{@{}l Y Y Y@{}}
\toprule
 & Idea & Closest prior work & Status here \\
\midrule
1 & SC-style carry, no BPTT ($K{=}1,n{=}0$) & RIN latent self-conditioning; Loopholing & tested; weakest \\
2 & Re-use computation $\Rightarrow$ fewer layers & implicit ML force fields (warm-started fixed point, 2--5$\times$) & GEOM 2-layer result \\
3 & BPTT: latent optimised for future steps & Relay (TBPTT, horizon 2); MetaState & $n{\ge}1$ helps \\
4 & Train like inference: update $t$, update $x_t$, long trajectories & MetaState ($K$-step unroll along the reveal schedule); FPF (on-policy carry) & \textbf{untested} $\rightarrow$ \S\ref{sec:main} \\
5 & What to backprop through (latent only, or also the integrator) & --- & untested \\
6 & Loss at every step & TRM deep supervision; MetaState & \textbf{untested} $\rightarrow$ \S\ref{sec:main} \\
7 & Same-$(x_t,t)$ recursion with BPTT; hierarchical inference & TRM; DEQ / Jacobian-free backprop; FPF & current training; multi-level inference was null \\
8 & Learning-rate scaling $1/n$ & --- & see note in \S\ref{sec:cost} \\
9 & Tiny models that need many passes; loss only after warm-up & TRM & the context floor $m_0$, \S\ref{sec:main} \\
10 & Looped layers inside one pass & Looped Diffusion LMs & control, \S\ref{sec:secondary} \\
11 & Iterate to a fixed point, adaptive depth & Fixed-point flows; DEQuiformer & later (inference) \\
12 & Residual-style gradient path between steps & --- & \code{lin\_res} is this \\
\bottomrule
\end{tabularx}
\end{center}

% ======================================================================
\section{A working picture: tracking a moving fixed point}\label{sec:picture}

Write one network pass as $(\hat y, z') = F_\theta(x,\sigma,z)$. For fixed $(x,\sigma)$,
iterating $z \leftarrow F_\theta(x,\sigma,z)$ heads to a fixed point $z^\star(x,\sigma)$.

\begin{itemize}
  \item \textbf{Current training} is the implicit-model recipe: half the time $K$
        untracked iterations to get near $z^\star$, then always $n$ tracked ones, loss
        on the readout of the last.
  \item \textbf{Inference} does one iteration per $\sigma$-step, warm-started from the
        neighbouring step. It tracks a slowly moving fixed point, always a little
        behind: lag $\approx$ drift per step$/(1-\rho)$ for contraction rate $\rho$.
\end{itemize}

This one picture accounts for most of what has been observed:
\begin{itemize}
  \item \emph{$K{>}1$ helps.} Carries must come from the model's own multi-pass
        dynamics. This is exactly the finding of Fixed-Point Forcing: one-pass carries
        give exposure bias and the recursion settles at wrong fixed points.
  \item \emph{Multi-level inference was null.} The carried latent is already near
        $z^\star$; an extra same-$\sigma$ iteration is worth about as much as a smaller
        $\Delta\sigma$.
  \item \emph{More sampling steps help LR more than vanilla} (thesis table). Smaller
        drift per step, smaller lag.
  \item \emph{The latent adds computation, not information.} The ideal denoiser is a
        function of $(x_t,\sigma)$ alone; the carry is a warm start (the fixed-point-flows
        paper says the same: the coupling to the carried state is ``a byproduct of warm
        starts''). That is why LR makes 2 layers behave like 4, and it constrains
        which training couplings are legitimate (\S\ref{sec:leak}).
\end{itemize}

\textbf{The gap this exposes.} Training asks one question: after 6 or 12 passes on the
\emph{same} input, is the readout good? Inference asks a different one: given a latent
made one step ago at a \emph{noisier} input, is the readout good after a single pass?
And seven passes sit in the autograd graph at every step while six of their
predictions are thrown away.

% ======================================================================
\section{What actually changes between two inference passes}\label{sec:gap}

Current training holds all of these fixed. From \code{\_generate}, between the pass
that produces the latent and the pass that consumes it:

\begin{center}
\begin{tabularx}{\textwidth}{@{}l Y l@{}}
\toprule
Change & What happens at inference & In training \\
\midrule
(A) $\sigma$ & decreases one schedule step; \code{c\_noise}$=\log\sigma/4$ shifts & fixed \\
(B) coordinates & $x \mathrel{+}= \Delta\,(\xh - x)/\sigma$: drift \emph{toward the model's own prediction} & fixed \\
(B$'$) churn & fresh noise added when $\sigma\in[0.03,3]$ & absent \\
(C) discrete sites & a fraction unmasked by \emph{sampling} from the predicted distribution; irreversible & none ever revealed \\
(D) \code{cond\_batch} & refreshed from the latest, latent-informed prediction & frozen after one cold pass \\
\bottomrule
\end{tabularx}
\end{center}

\paragraph{(A) is staleness.} The walk in \S\ref{sec:main} fixes it, leak-free and at
no extra cost.

\paragraph{(B) is exposure bias, and LR amplifies it.} At inference $x_{i+1}$ has
already been dragged toward $\xh^{(i)}$ --- the very thing the latent encodes. If that
prediction was bad, the state now sits off the noising manifold \emph{in the direction
of the model's own error}. Teacher-forced training (including the walk below) never
produces such a state: $x^{(i+1)}$ is a clean noising of $x_0$ however badly the model
did. Every diffusion model has this mismatch; LR makes it sharper, because a past
mistake is carried forward explicitly in the latent instead of only implicitly through
$x_t$ where the schedule contracts it. Plausibly why LR-without-SC is unstable on GEOM.
Note (B$'$) partly rescues teacher-forcing in $\sigma\in[0.03,3]$, since inference
re-noises there too.

\paragraph{(C) is the largest, and the cheapest to fix.} \code{\_mask\_sampling\_step}
commits sites by sampling from the predicted distribution, and in the mask formulation
they never revert. So at inference some sites have just collapsed from a distribution
to a hard value; that value may contradict what the latent believed; and the model must
update its remaining beliefs conditioned on it. Current training reveals \emph{nothing}
across the whole chain, so the model never experiences a commitment event. A
teacher-forced walk would reveal from ground truth --- always correct, so it would
never see a \emph{wrong} commitment.

This matters most because validity is purely a discrete-head metric
(\code{mol is None} $\Leftrightarrow$ sanitisation fails on the predicted
type/bond/charge graph), and validity is where LR is winning on GEOM.

\textbf{The fix: reveal from the model's own samples, not from $x_0$.} The loss
machinery already supports it --- \code{\_type\_loss} and \code{\_bond\_loss} restrict
to sites masked in \code{interpolated}:
\begin{lstlisting}
masked_types = argmax(interpolated["atomics"]) == type_mask_index
type_loss = type_loss * masked_types
\end{lstlisting}
so self-revealing automatically moves the loss onto the still-masked sites and leaves
the target valid. Coordinates regress to $x_0$ as before. Nothing else changes.

\paragraph{Order of attack.} (A) then (C) then (B). (B) last because it is the one that
breaks the regression target: regressing to $x_0$ from a state that is no longer a
noising of $x_0$. The standard mitigation is scheduled sampling --- on-policy with
probability $\rho$, ramped during training.

\paragraph{Measure the gap before changing anything.} No retraining. On validation data
at fixed $\sigma$, compare the $\xh$ loss under three protocols:
\begin{enumerate}[label=P\arabic*.]
  \item fixed $x_t$, 6 passes --- the training configuration;
  \item cold start, 1 pass;
  \item latent carried from a genuine previous inference step (run the real loop from a
        noised validation molecule).
\end{enumerate}
P1 $\ll$ P3 is the size of the gap, in loss units. P2 vs P3 says how much the carried
latent is actually buying in the deployment regime: if P3 $\approx$ P2 the latent is not
surviving the transition and (A)+(C) are urgent; if P3 $\ll$ P2 it is surviving and the
gap is smaller than feared.

% ======================================================================
\section{Main proposal: walk $\sigma$ along the tracked chain and supervise it}\label{sec:main}

Keep $K$, $n$, $p$, the pass count, the graph and the single optimizer step. Change
only what the tracked passes see and what they are asked.

\begin{algorithm}[h]
\caption{Proposed training step}\label{alg:traj}
\begin{algorithmic}[1]
  \State sample $\sigma_0$; take consecutive schedule points $\sigma_0>\sigma_1>\dots>\sigma_n$
  \State build a forward-Markov path, clean end first:
         $x^{(n)} = x_0+\sigma_n\varepsilon$,\;
         $x^{(i-1)} = x^{(i)} + \sqrt{\sigma_{i-1}^2-\sigma_i^2}\;\xi_i$;\; masks nested
  \State carry $c \gets \mathbf 0$, \quad $c = (\zs,\ze,\xh,p_{\text{type}},p_{\text{bond}})$
  \State \textbf{if} $u<p$: \textbf{repeat} $K$ times: $c \gets \sg\,F_\theta(x^{(0)},\sigma_0,c)$ \Comment{untracked warm-up, as today}
  \State \textbf{for} $i=0..n$: $(\hat y_i, c) \gets F_\theta(x^{(i)},\sigma_i,c)$ \Comment{tracked, as today, but stepping down $\sigma$}
  \State $S \gets \{\, i : \text{passes before pass } i \;\ge\; m_0 \,\}$ \Comment{context floor}
  \State \textbf{return} $\tfrac{1}{|S|}\sum_{i\in S}\lambda(\sigma_i)\,\ell(\hat y_i,x_0)$
\end{algorithmic}
\end{algorithm}

Three changes against Algorithm~\ref{alg:current}:
\begin{enumerate}[label=(\alph*)]
  \item the tracked passes step down the sampling schedule instead of repeating $(x_t,\sigma)$;
  \item every tracked pass with at least $m_0$ passes behind it gets a loss;
  \item the SC inputs are the previous pass's outputs (detached), not a frozen cold pass.
\end{enumerate}
Default $m_0 = n = 6$: on the cold branch only the last pass qualifies, which is today's
loss; on the warm branch all seven do. Four supervised predictions per step on average,
for the same compute. Adding self-sampled reveals (\S\ref{sec:gap}) on top is the
natural second increment.

\subsection{Why this should work}

\begin{enumerate}
  \item \textbf{The supervision is free.} The seven tracked passes and their readouts
        are already computed and already in the graph. (b) adds loss evaluations and
        nothing else.
  \item \textbf{Shorter gradient paths.} Each pass gets a direct gradient instead of
        one that has travelled through up to six further passes. Deep supervision is
        what keeps weight-tied recursion trainable in TRM, and looped models report
        instability growing with loop count. It may be the fix for LR-without-SC being
        unstable on GEOM.
  \item \textbf{Every supervised pass is in the deployment state}: a latent made at
        the noisier neighbour, one pass, readout. The fixed-$\sigma$ chain trains
        refinement in place, which inference never does, and which bought nothing when
        it was tried (multi-level inference).
  \item \textbf{BPTT says what it was meant to say.} The gradient from pass $i{+}1$'s
        loss into pass $i$ is literally ``make a latent the next, less noisy step can
        use''. Depth 6 is kept, so nothing is given up if deep BPTT is what small
        models need (the thesis held $(n{+}1)\times$layers constant).
  \item \textbf{Self-conditioning moves inside the chain.} Today $c$ comes from one
        latent-free pass and is frozen for up to 13 passes, while at inference it is
        refreshed every step by a latent-informed pass. With the carry
        $c=(\zs,\ze,\xh,p_{\text{type}},p_{\text{bond}})$ training matches inference,
        and the separate SC pass and its coin flip disappear.
  \item \textbf{It has precedent on the half of VEDA that decides validity.} Relay
        trains a carried per-token latent in masked diffusion with truncated BPTT at
        horizon 2; MetaState unrolls $K$ denoising steps along the reveal schedule with
        a loss at every step and BPTT, and reports that both the recurrence and the
        BPTT matter. Atom types and bonds in VEDA are masked discrete diffusion.
\end{enumerate}

\subsection{The path must be forward-Markov, or the latent leaks $x_0$}\label{sec:leak}

How consecutive inputs $x^{(i-1)},x^{(i)}$ are coupled matters more than it looks.
\begin{itemize}
  \item \textbf{Shared noise}, $x^{(i)}=x_0+\sigma_i\varepsilon$: two points fix the
        line, so a latent that remembers the previous geometry gives $x_0$ exactly. At
        inference the same extrapolation only returns the previous prediction. The
        model would learn to trust its own last output completely --- the same failure
        as training SC on the true $x_0$.
  \item \textbf{Independent noise} per step: history is several independent
        observations of $x_0$. The model learns to treat the carry as fresh evidence;
        at inference consecutive states are strongly correlated, so it becomes
        over-confident.
  \item \textbf{Forward-Markov}, $x^{(i-1)}=x^{(i)}+\text{fresh noise}$: $x_0$ is
        independent of the past given the present. History carries computation only,
        which is the regime the fixed-point picture says is the true one. This is the
        joint law of the exact reverse SDE. For the discrete channels it means nested
        masks (a site masked at $\sigma_i$ is masked at every noisier level): draw one
        uniform per site and mask iff $u<r(\sigma)$.
\end{itemize}
The current fixed-$x_t$ chain is the $\Delta\sigma=0$ case of this, which is why it is
leak-free. Forward-Markov neighbours are further apart than consecutive sampler states
(for a geometric 100-step schedule, about $0.5\sigma$ of fresh noise against about
$0.1\sigma$ of deterministic motion), so training sees \emph{more} lag than inference:
conservative in the right direction. \code{\_interpolate\_with\_noise} already moves a
state to a noisier level; check that its categorical branch re-masks with the
mask-rate difference.

\subsection{Your concern: per-step losses and the small-model regime}

Today every supervised prediction has 6 or 12 passes behind it; the model is never
asked to be good with less. A loss on earlier passes would ask exactly that, which is
your worry, and it is a fair one. The context floor $m_0$ keeps the property: at
$m_0=n$ no supervised prediction has less context than today's minimum.

$m_0$ is then the dial for your hypothesis. Lowering it ($m_0=3$, then $0$) asks the
model to be good earlier. If final-state-only supervision is what lets tiny models
work, validity should fall as $m_0$ falls.

\subsection{Cost and learning rate}\label{sec:cost}

\begin{center}
\begin{tabular}{lcc}
\toprule
per optimizer step & current & proposed \\
\midrule
forward passes (expected) & 10.5 & 10.5 \\
tracked passes            & 7 & 7 \\
supervised predictions    & 1 & 4 \;($m_0{=}6$),\; 5.5 \;($m_0{=}3$) \\
\bottomrule
\end{tabular}
\end{center}

A step costs the same by construction; it teaches more. Making a step cheaper is a
separate lever (\S\ref{sec:secondary}, items 1--2).

On the $1/n$ question: Adam divides out the gradient scale, so summing gradients over
tied passes does not by itself call for a smaller learning rate; the clipping threshold
is the thing that interacts. The loss above is a mean, so its scale stays put.

Two details. The supervised $\sigma_i$ sit up to $n$ schedule steps below $\sigma_0$
(about a factor 2 for six steps of a geometric 100-step schedule), so sample $\sigma_0$
half a walk higher to keep them centred on the training distribution. And the module
has to build the path itself from $x_0$ and the prior sample, with one uniform per
site for the masks.

\subsection{Main experiment}

GEOM-Drugs, 2 layers (QM9 validity has no headroom). A $2\times2$:

\begin{center}
\begin{tabular}{lcc}
\toprule
 & loss on last pass only & loss on tracked passes ($m_0{=}6$) \\
\midrule
fixed $\sigma$   & A (exists)  & B$'$ \\
$\sigma$-walk    & B$''$       & \textbf{B} \\
\bottomrule
\end{tabular}
\end{center}

B$'$ is a few lines and can run today. All four cost the same per step, so compare at
equal steps on validity, fc-validity and opt-energy-validity (fixed denominators), then
molecule stability and opt-rmsd.

\emph{Predictions.} B$'$ $>$ A in speed and stability. B$''$ $>$ A in sample quality
even where its train loss is no better. B is best, and the two walk arms lose less when
sampled with 50 or 25 steps. \emph{If instead} B$'$ $<$ A, your hypothesis holds:
keep the last-pass loss, and the question reduces to B$''$ against A.

% ======================================================================
\section{The coordinate channel}\label{sec:coords}

\subsection{Why carrying $v$ did nothing}

\begin{enumerate}
  \item The carry adds computation, not information. So the question is which
        geometric computation is expensive for two layers and reusable.
  \item The trunk reads geometry only through invariants: Gram products
        $\langle x_i,x_j\rangle$ in \code{EdgeMessages}, lengths in \code{LengthsMLP},
        distances in \code{BondRefine} (output end only). The vector stream itself is
        only mixed linearly and read by \code{coord\_head}.
  \item A linear mix of $v_{\text{prev}}$ into the coord sets is therefore visible
        downstream only through the Gram of the \emph{mixed} sets: current, cross and
        previous terms tangled in the same 64 channels. The easy useful solution is
        rank one --- pick out the \code{coord\_head} direction, i.e.\ $\xh$ --- which is
        just coordinate self-conditioning. Hence \code{v\_lin\_res} $\approx$ SC.
  \item $\ze$ already \emph{is} the geometry carry. \code{BondRefine} builds it from
        $\|\xhat{i}-\xhat{j}\|^2$, $\langle\xhat{i},\xhat{j}\rangle$ and the node
        pairs. The coordinate information crosses the step in the only form the trunk
        reads.
  \item What is poorly served is \emph{local distance}. Zero-centred dot products
        encode a bond length as a small difference of large numbers (worse for big
        molecules), there is no RBF expansion anywhere, and true distances appear only
        at the output. In the thesis, RBF distances from the prediction on the edge
        route were the largest single contributor to self-conditioning.
\end{enumerate}

\subsection{Proposal: geometry crosses the step as distances}

Carry $\xh$ (it is free: the pass that produced $z$ produced it), and use it twice:
\begin{align*}
  \text{equivariant:}\quad & \code{coord\_proj}\big(\,[\,c_{\text{in}}x_t,\ \xh\,]\,\big) && \text{(as SC does today, now refreshed in-chain)}\\
  \text{invariant:}\quad & e \leftarrow e + W\,\mathrm{RBF}\!\big(\|\xhat{i}-\xhat{j}\|\big) && \text{at \code{e1}, } W \text{ zero-init, } 32{\times}128 \approx 4\text{k params}
\end{align*}
$\xh$ is in data units at every $\sigma$, so the RBF grid is fixed (0 to $\sim$10\,\AA).
Zero at a cold start. Detach $\xh$ first; let gradients through it later.

If you want the full 64-set latent rather than its readout: feed the self-Gram
$\langle v_{\text{prev},s,i}, v_{\text{prev},s,j}\rangle$ (64 invariants per pair) into
\code{e1} through a zero-init linear map, 8k parameters. I expect this to be mostly
redundant with $\ze$; it is the honest test of whether the other 63 directions hold
anything.

% ======================================================================
\section{Secondary experiments}\label{sec:secondary}

Separate from the main experiment; ordered by value for cost.

\begin{enumerate}
  \item \textbf{$K6/n0$ vs $K6/n1$ vs $K6/n6$} on the current scheme. $K6/n0$ is
        Fixed-Point Forcing, which reports needing no BPTT at all. Tracked passes are
        the expensive ones (7 of them, every step), so if $n{=}1$ matches $n{=}6$ the
        method gets much cheaper for free.
  \item \textbf{Persist and truncate.} Keep the carry alive across optimizer steps on
        the same batch, backprop through windows of $w$ passes, loss on each: one
        supervised prediction per forward pass and $w$ graphs in memory instead of 7.
        This is how TRM and ordinary RNN training amortise the hidden state. Worth
        doing once the main experiment shows per-pass losses are safe.
  \item \textbf{Self-sampled discrete reveals} (\S\ref{sec:gap}). Small change on top
        of the walk; the loss masking already supports it; aimed straight at validity.
  \item \textbf{$m_0$ sweep} ($3$, $0$) on arm B. The direct test of whether
        final-state-only supervision is what enables small models.
  \item \textbf{SC channel ablation} in VEDA-S: coordinates only, discrete only. Tells
        whether \emph{any} geometric carry matters here, before more is built.
  \item \textbf{RBF distance injection without LR} (from $\xh$ via SC). Tests point 5 of
        \S\ref{sec:coords} on its own; could help every few-layer model.
  \item \textbf{Looped-layer control.} A 3-layer model with the core layer applied $T$
        times per pass: same parameters as 3 layers, depth per pass $T{+}2$. The natural
        parameter-matched competitor to cross-step recursion.
  \item \textbf{Fewer sampling steps.} Sample the main-experiment arms at 100/50/25
        steps; in the walk arms randomise the stride over $\{1,2,4\}$ schedule points.
        If LR models degrade more slowly, that is an inference-cost result.
  \item \textbf{Oracle latent.} At inference, replace the carry with the latent from a
        pass on the ground-truth molecule. No gain $\Rightarrow$ the mechanism cannot
        use the channel. Large gain $\Rightarrow$ the model cannot yet produce good
        latents; spend on training.
  \item \textbf{Log gates and latent norms} ($\|W\|$ of each transfer; $\|\zs\|,\|\ze\|$
        per sampling step).
  \item \textbf{On-policy $x$ with scheduled sampling.} Advance $x$ with the integrator
        from the model's own prediction instead of the forward-Markov path, with
        probability $\rho$ ramped during training. Matches inference exactly but
        regresses towards $x_0$ from states that are no longer its noised versions.
\end{enumerate}

% ======================================================================
\section{Literature notes}

From abstracts and secondary summaries; arXiv was not reachable from this session, so
none of these was read in full. Details marked ``reported'' should be checked.

\begin{itemize}
  \item \textbf{Flow Reasoning Models} (\href{https://arxiv.org/abs/2606.29150}{2606.29150}).
        Conventional SC degrades with recurrent depth: exposure bias between one-pass
        training carries and recursively generated inference states; loss rises with
        depth as predictions settle at wrong fixed points. \emph{Fixed-Point Forcing}
        forms the detached carry by repeated application of the denoiser; reported to
        need no BPTT. This is the present scheme with $K{>}1$, $n{=}0$.
  \item \textbf{Fixed-point flows} (\href{https://arxiv.org/abs/2607.00714}{2607.00714}).
        SC in flow language models is a fixed-point iteration on the denoising estimate;
        under contractivity the fixed point is unique, and replacing the carry by its
        fixed point leaves an ordinary flow in $x_t$ alone. Supports
        \S\ref{sec:picture}.
  \item \textbf{Implicit ML force fields} (\href{https://arxiv.org/abs/2607.29158}{2607.29158});
        \textbf{DEQuiformer} (\href{https://arxiv.org/abs/2509.08734}{2509.08734}).
        Equivariant GNNs as fixed-point models; the previous MD step's fixed point
        warm-starts the next. Reported 2--5$\times$ compute and memory; solver steps
        from 5--6 to about 3. The same structure as LR at inference, in a
        supervised setting.
  \item \textbf{Relay} (\href{https://arxiv.org/abs/2605.22967}{2605.22967}). Learned
        per-token latent carried across denoising steps of a masked diffusion model;
        truncated BPTT at horizon 2.
  \item \textbf{Loopholing}. Deterministic latent pathway across steps, trained
        SC-style without unrolling.
  \item \textbf{MetaState} (\href{https://arxiv.org/abs/2603.01331}{2603.01331}).
        Memory slots carried across denoising steps of a frozen discrete diffusion LM;
        $K$-step unrolling along a sampled reveal schedule, loss at each step, BPTT;
        ablations reported to show recurrence and BPTT both matter.
  \item \textbf{Looped Diffusion Language Models} (\href{https://arxiv.org/abs/2605.26106}{2605.26106}).
        Looping the early-middle layers of a masked diffusion transformer improves
        training efficiency and quality. (2606.18022 and 2605.31215 could not be
        identified by search.)
  \item \textbf{RIN} (\href{https://arxiv.org/abs/2212.11972}{2212.11972}). Latent
        self-conditioning: previous-step latents warm-start the current step, trained
        with stop-gradient at the same $t$, injected through a zero-initialised
        normalisation.
  \item \textbf{TRM}. Deep supervision with a detached carry over repeated steps on the
        same sample; gradient through the last recursion only.
\end{itemize}

% ======================================================================
\appendix
\section{Metric definitions}\label{app:metrics}

\begin{lstlisting}
# val-validity                 == count(mol is not None) / N
#   mol_from_atoms(sanitise=True) returns None on failure; mol_is_valid then
#   re-runs the same idempotent SanitizeMol, so it always passes.
# val-fc-validity              same, plus len(GetMolFrags(mol)) == 1

# val-opt-energy-validity      denominator = ALL samples
n_ok = 0
for mol in all_generated:
    if mol is None: continue
    opt = optimise_mol(mol)                 # None only if MMFF raises
    if opt is None: continue
    if not isfinite(calc_energy(opt)): continue
    n_ok += 1
return n_ok / len(all_generated)

# val-energy                   no optimisation; raw pose; valid mols only
es = [calc_energy(m) for m in all_generated if m is not None]
return mean([e for e in es if isfinite(e)])

# val-opt-rmsd                 built AND optimisable mols only
rmsds = [rmsd(m, optimise_mol(m)) for m in built_and_optimisable]
return mean(rmsds)                          # centre, Kabsch, sqrt(mean sq dist)

# val-strain(-per-atom)        energy stored in the pose
diffs = [calc_energy(m) - calc_energy(optimise_mol(m))
         for m in built_and_optimisable]
return mean(diffs)

# val-novelty                  valid SMILES only
return mean([canonical_smiles(m) not in train_smiles for m in valid])
\end{lstlisting}

\begin{itemize}
  \item Validity is chemical, not geometric: bonds come from the edge head, valence is
        checked against the predicted charge, aromatic is a predicted class.
  \item Only validity, fc-validity and opt-energy-validity have a model-independent
        denominator. Energy, opt-rmsd and strain average over a subset that grows with
        model quality.
  \item Strain is conformational: energy relative to the molecule's own relaxed pose.
        It is dominated by the worst contact; opt-rmsd averages over atoms.
  \item \code{optimise\_mol} ignores the convergence flag; \code{AverageStrainEnergy}
        does not filter \code{inf}/\code{nan}.
\end{itemize}

\end{document}
